\reportchapter{Conception et développement de la partie logicielle}
\label{chap:partie_logicielle}

\section{Introduction}
\label{sec:ch5_introduction}

Les chapitres précédents ont établi le modèle du procédé puis le BMS sous test. La partie logicielle du projet répond à un autre besoin : fournir à l'ingénieur BMS un outil capable de lancer les validations, collecter les résultats, analyser les modèles et produire un rapport exploitable sans manipuler séparément les outils techniques sous-jacents.

\noindent
Le logiciel développé n'est donc pas un composant embarqué du BMS. Il constitue une couche d'assistance d'ingénierie placée autour des modèles. Son rôle est de rendre la chaîne de validation plus reproductible, plus lisible et plus facile à auditer. La conception suit une logique simple : chaque ensemble possède une responsabilité claire, chaque résultat important est conservé sous une forme structurée, et les ensembles communiquent par des données inspectables plutôt que par des états implicites.

\section{Architecture logicielle générale}
\label{sec:ch5_architecture}

L'architecture logicielle repose sur cinq ensembles principaux. Le validateur exécute les campagnes de validation modèle, code hôte et cible. L'analyseur interprète les résultats et extrait des métriques statiques des modèles. Le générateur de rapports consolide les données et produit des synthèses lisibles. L'orchestrateur enchaîne les phases. L'interface MATLAB expose l'ensemble à l'utilisateur final.

\begin{figure}[H]
  \centering
  \fbox{\parbox{0.88\textwidth}{\centering\vspace{2.2cm}
  \textit{Emplacement réservé à une figure d'architecture logicielle.}\\[0.4em]
  \small La figure doit montrer le validateur, l'analyseur, le générateur de rapports, les orchestrateurs et l'interface MATLAB, avec les données structurées échangées.
  \vspace{2.2cm}}}
  \captionsetup{width=0.84\textwidth, skip=15pt}
  \caption{Architecture logicielle de l'assistant d'ingénierie BMS.}
  \label{fig:ch5_architecture_logicielle}
\end{figure}

\subsection{Principe de séparation des responsabilités}
\label{subsec:ch5_separation}

Le découpage retenu évite de concentrer toute la logique dans une application unique. Les traitements liés aux modèles restent proches de l'environnement de simulation et des campagnes MIL, SIL et PIL. Les traitements d'analyse prennent en charge la consolidation, l'application de règles et la génération de rapports. Cette séparation permet de tester chaque partie indépendamment et de remplacer une brique sans modifier toute la chaîne.

\begin{table}[H]
  \centering
  \small
  \begin{tabularx}{\textwidth}{|l|X|X|}
    \hline
    \textbf{Ensemble} & \textbf{Moyens principaux} & \textbf{Responsabilité} \\
    \hline
    Validateur & Fonctions de validation proches des modèles & Exécuter les scénarios, comparer les signaux, produire les verdicts et les figures. \\
    \hline
    Analyseur & Collecteurs de faits et moteur de règles & Extraire les faits des modèles, appliquer les règles et produire un diagnostic hiérarchisé. \\
    \hline
    Générateur de rapports & Chaîne de consolidation et de synthèse & Consolider les données, réduire le volume utile et générer les rapports audités. \\
    \hline
    Orchestrateurs & Points d'entrée automatisé et interactif & Enchaîner les phases en mode automatisé ou depuis l'interface. \\
    \hline
    Interface & Application graphique MATLAB & Fournir les vues de pilotage, de consultation et d'inspection au même ingénieur BMS. \\
    \hline
  \end{tabularx}
  \captionsetup{width=0.9\textwidth, skip=15pt}
  \caption{Répartition des responsabilités entre les ensembles logiciels.}
  \label{tab:ch5_responsabilites}
\end{table}

\subsection{Familles d'exigences logicielles}
\label{subsec:ch5_exigences_logicielles}

Les exigences logicielles sont regroupées par famille fonctionnelle. Le texte du chapitre conserve cette logique de regroupement pour ne pas alourdir la lecture avec une liste d'identifiants trop longue. Le catalogue complet des exigences modèle et logiciel est reporté en annexe, où le lecteur peut inspecter les identifiants et le contenu compacté.

\begin{table}[H]
  \centering
  \small
  \begin{tabularx}{\textwidth}{|l|c|X|}
    \hline
    \textbf{Famille} & \textbf{Nombre} & \textbf{Périmètre} \\
    \hline
    Orchestration automatisée & 3 & Exécution contrôlée des phases depuis un point d'entrée automatisé. \\
    \hline
    Génération de rapports & 7 & Consolidation, réduction, génération de synthèse et audit des références. \\
    \hline
    Analyse & 7 & Construction du contexte d'analyse, moteur de règles et hiérarchisation des constats. \\
    \hline
    Orchestration interactive & 5 & Pilotage des campagnes depuis l'interface. \\
    \hline
    Validation & 5 & Aides de validation, extraction de signaux, scénarios, figures et export structuré. \\
    \hline
    Collecte de faits & 2 & Collecte des métriques statiques des modèles. \\
    \hline
    \textbf{Total} & \textbf{29} & Exigences logicielles documentées et reliées aux tests. \\
    \hline
  \end{tabularx}
  \captionsetup{width=0.88\textwidth, skip=15pt}
  \caption{Regroupement des exigences logicielles par famille.}
  \label{tab:ch5_familles_req_sw}
\end{table}

\section{Module de validation}
\label{sec:ch5_validateur}

Le validateur est la base expérimentale de l'outil. Il transforme les modèles Simulink et les scénarios de test en rapports structurés. Son architecture reprend les niveaux de validation utilisés pour le BMS : validation du procédé, MIL, SIL et PIL. Les fonctions transverses assurent la préparation des entrées, la récupération des signaux, la comparaison des résultats et l'export des rapports.

\subsection{Organisation des campagnes}
\label{subsec:ch5_campagnes}

La validation du procédé reste séparée du pipeline BMS. Elle évalue la fidélité du modèle physique sur les trajets réels. La validation BMS suit ensuite une progression MIL, SIL et PIL. Le MIL vérifie les fonctions sur modèle, le SIL compare le comportement du code généré sur hôte avec la référence Simulink, et le PIL vérifie l'exécution sur la cible STM32.

\noindent
Cette organisation évite de mélanger deux questions différentes. Le procédé doit d'abord être suffisamment représentatif pour servir d'environnement. Le BMS est ensuite évalué dans cet environnement selon des niveaux de confiance croissants.

\begin{table}[H]
  \centering
  \small
  \begin{tabularx}{\textwidth}{|l|X|X|}
    \hline
    \textbf{Niveau} & \textbf{Objectif} & \textbf{Sorties principales} \\
    \hline
    Procédé & Comparer tension pack, SoC et température aux mesures BMW i3. & Métriques de convergence et figures par trajet. \\
    \hline
    MIL & Vérifier les fonctions BMS et prédicteur avant génération de code. & Rapport structuré, figures de scénario et données de couverture. \\
    \hline
    SIL & Vérifier l'équivalence entre modèle et code généré sur hôte. & Rapport structuré et figures d'équivalence. \\
    \hline
    PIL & Vérifier la cohérence numérique et le temps d'exécution sur cible embarquée. & Rapport structuré, mesures de temps et figures associées. \\
    \hline
  \end{tabularx}
  \captionsetup{width=0.9\textwidth, skip=15pt}
  \caption{Rôle des niveaux de validation pilotés par le logiciel.}
  \label{tab:ch5_niveaux_validation}
\end{table}

\subsection{Rapports et figures}
\label{subsec:ch5_rapports_validateur}

Les rapports du validateur utilisent une structure volontairement simple. Chaque test porte un identifiant d'exigence, un nom, un statut et une liste de vérifications élémentaires. Les figures générées complètent ces rapports en montrant les signaux d'entrée, les réponses attendues, les seuils et les verdicts. Ce choix facilite l'audit : une synthèse peut être lue rapidement, mais les détails restent disponibles pour remonter jusqu'au signal qui justifie le verdict.

\section{Module d'analyse}
\label{sec:ch5_analyseur}

L'analyseur ajoute une couche d'interprétation au-dessus des résultats bruts. Il ne modifie pas les modèles et ne corrige pas automatiquement les défauts. Il collecte des faits, applique des règles et produit des constats classés par sévérité. Cette approche correspond mieux au rôle de l'outil : assister l'ingénieur BMS dans son diagnostic, tout en lui laissant la décision finale.

\subsection{Collecte des faits}
\label{subsec:ch5_collecte_faits}

Les collecteurs inspectent les modèles du maître, des esclaves et du prédicteur. Ils extraient des informations sur la structure des modèles, les automates, les solveurs, les signaux, les blocs de calcul, les données et la génération de code. Le résultat est regroupé dans un rapport de faits qui sert de base stable au moteur d'analyse.

\noindent
Le collecteur a été conçu pour rester non intrusif. Il lit les modèles, filtre les faux positifs liés aux blocs de bibliothèque et agrège les défauts de nommage par parent afin d'éviter une accumulation de messages peu exploitables.

\subsection{Moteur de règles}
\label{subsec:ch5_moteur_regles}

Le moteur d'analyse charge les faits, les rapports de validation, les exigences modèle et une base de connaissances déclarative. Les règles couvrent notamment la couverture structurelle, l'architecture du modèle, les conventions de modélisation et les résultats MIL, SIL et PIL. Chaque constat contient une sévérité, une explication, une recommandation et, lorsque cela est possible, un lien de localisation vers l'élément concerné dans Simulink.

\noindent
La base de règles centralise les seuils par composant. Le BMS maître est traité comme le composant le plus critique, le BMS esclave comme un composant de surveillance locale, et le prédicteur comme une fonction informative. Cette séparation évite de donner le même poids à un défaut sur la chaîne de sûreté déterministe et à un défaut sur une fonction d'aide à l'anticipation.

\section{Module de génération de rapports}
\label{sec:ch5_generateur}

Le générateur de rapports transforme les données techniques en documents lisibles. Il est organisé autour de trois idées : conserver une version complète non modifiée des données, construire une version réduite pour la génération de texte, puis auditer toutes les références présentes dans le résultat final.

\subsection{Consolidation des données}
\label{subsec:ch5_consolidation}

Le rapport complet conserve les résultats MIL, SIL, PIL et analyse sous une structure unique. Il n'est pas réécrit ni simplifié. Une vue plus légère est ensuite construite pour la génération automatique. Chaque entrée importante conserve un lien de traçabilité vers sa source dans le rapport complet.

\noindent
Ce mécanisme est essentiel pour éviter une synthèse non vérifiable. Le rapport généré peut être agréable à lire, mais toute affirmation technique doit rester reliée à un emplacement précis dans les données brutes.

\subsection{Génération de synthèse en deux parties}
\label{subsec:ch5_generation_html}

La génération de synthèse est divisée en deux parties. La première partie présente les tableaux de validation. La seconde partie présente l'analyse d'ingénierie et le verdict. Cette séparation réduit la taille de chaque requête envoyée au modèle de langage et rend le résultat plus stable. Le générateur peut également fonctionner en mode de simulation, dans lequel seules les données structurées sont produites. Ce mode reste utile pour les tests et pour les situations où aucun service de génération textuelle n'est disponible.

\subsection{Audit des références}
\label{subsec:ch5_audit_refs}

Après génération, les références citées dans la synthèse sont vérifiées. Elles doivent résoudre vers une entrée réelle du rapport complet. Cette étape empêche qu'une référence inventée soit publiée comme preuve.

\section{Orchestration du pipeline}
\label{sec:ch5_orchestrateur}

Deux orchestrateurs coexistent. Le premier est destiné aux exécutions automatisées. Le second est utilisé par l'application graphique. Leur coexistence répond à deux usages différents : automatiser une campagne depuis un environnement de validation, ou piloter la même logique depuis l'environnement de travail de l'ingénieur BMS.

\subsection{Orchestrateur automatisé}
\label{subsec:ch5_orch_python}

L'orchestrateur automatisé construit la séquence d'exécution, enchaîne les phases demandées, puis lance l'analyseur et le générateur de rapports. Il accepte des options de saut de phase, ce qui permet de réutiliser des résultats existants ou de contourner le PIL lorsque la carte cible n'est pas connectée. Cette variante est adaptée aux campagnes reproductibles et aux usages proches de l'intégration continue.

\subsection{Orchestrateur MATLAB}
\label{subsec:ch5_orch_matlab}

L'orchestrateur interactif fournit le pilotage intégré à l'interface. Il accepte les niveaux de validation sélectionnés, les filtres MIL, SIL et PIL, l'option de réutilisation du code généré, ainsi que le suivi du journal et de la progression. La couverture est collectée avec le MIL lorsqu'elle est pertinente, puis les faits, l'analyse et le rapport sont produits dans la continuité.

\noindent
Le catalogue de scénarios alimente les listes de l'interface. Les scénarios MIL sont regroupés par famille fonctionnelle, tandis que les scénarios SIL et PIL restent plus compacts. Cette organisation donne à l'utilisateur un contrôle suffisamment fin sans l'obliger à connaître l'organisation interne des tests.

\section{Interface utilisateur MATLAB}
\label{sec:ch5_ui}

L'interface est une application MATLAB intégrée à l'environnement déjà utilisé pour les modèles. Le public visé est un seul acteur : l'ingénieur BMS. Il peut toutefois utiliser l'outil dans deux postures différentes, une posture de développement lorsqu'il ouvre ou restaure les modèles, et une posture de validation lorsqu'il lance les campagnes et inspecte les rapports.

\subsection{Vues disponibles}
\label{subsec:ch5_vues}

L'application est organisée en cinq vues. La vue Dashboard donne l'état général et les actions rapides. La vue Validate permet de configurer et lancer les campagnes. La vue Reports liste les rapports générés et permet de les ouvrir, exporter ou supprimer. La vue Plots affiche les figures issues des validations. La vue Model donne accès aux modèles éditables, à la restauration de snapshots et à l'export du code PIL.

\begin{table}[H]
  \centering
  \small
  \begin{tabularx}{\textwidth}{|l|X|}
    \hline
    \textbf{Vue} & \textbf{Usage principal} \\
    \hline
    Dashboard & Visualiser l'état de l'outil, le dernier lancement et les actions rapides. \\
    \hline
    Validate & Choisir MIL, SIL, PIL, filtrer les scénarios, suivre la progression et consulter le journal. \\
    \hline
    Reports & Parcourir les rapports générés, les exporter ou les supprimer après confirmation. \\
    \hline
    Plots & Prévisualiser les figures produites par les tests. \\
    \hline
    Model & Ouvrir les modèles, restaurer les snapshots et exporter le code PIL disponible. \\
    \hline
  \end{tabularx}
  \captionsetup{width=0.86\textwidth, skip=15pt}
  \caption{Vues principales de l'application MATLAB.}
  \label{tab:ch5_vues_ui}
\end{table}

\subsection{Gestion de la progression}
\label{subsec:ch5_progression_ui}

L'exécution du pipeline reste dans le processus MATLAB principal, car l'initialisation des modèles dépend de l'environnement courant. L'interface compense cette contrainte par des mises à jour explicites de l'affichage. Les journaux et la barre de progression sont donc rafraîchis entre les étapes du pipeline, ce qui maintient une rétroaction visible pendant les campagnes longues.

\section{Vérification logicielle}
\label{sec:ch5_verification_logicielle}

La vérification logicielle combine des tests automatisés, des tests MATLAB ciblés et une recette manuelle de l'interface. Les tests automatisés couvrent l'analyseur, le générateur de rapports, la consolidation, les consignes de génération, le client de génération textuelle et l'orchestrateur automatisé. Les tests MATLAB disponibles ciblent notamment les scénarios, la liste des rapports et les aides de validation.

\noindent
La suite Python a été relancée pendant la préparation de cette version du rapport. Elle comporte 43 tests et se termine sans échec. La validation de l'interface reste volontairement manuelle, car elle porte sur des interactions graphiques, des confirmations utilisateur et la lisibilité des vues.

\begin{table}[H]
  \centering
  \small
  \begin{tabularx}{\textwidth}{|l|X|c|}
    \hline
    \textbf{Famille de vérification} & \textbf{Périmètre} & \textbf{Résultat disponible} \\
    \hline
    Tests automatisés & Analyseur, générateur de rapports, consolidation, consignes de génération, client de génération textuelle et orchestrateur automatisé. & 43 / 43 réussis \\
    \hline
    Tests MATLAB ciblés & Scénarios, liste des rapports et aides de validation. & Tests présents \\
    \hline
    Recette UI & Parcours des vues, lancement d'une campagne, consultation des rapports et restauration de modèle. & Recette manuelle \\
    \hline
  \end{tabularx}
  \captionsetup{width=0.88\textwidth, skip=15pt}
  \caption{Synthèse des vérifications logicielles disponibles.}
  \label{tab:ch5_verification_logicielle}
\end{table}

\section{Conclusion}
\label{sec:ch5_conclusion}

La partie logicielle réalisée fournit une chaîne cohérente autour des modèles BMS : elle lance les validations, structure les résultats, analyse les modèles, génère des rapports et expose l'ensemble dans une interface MATLAB. Le choix de données structurées inspectables donne une continuité entre exécution, analyse et justification. Le chapitre suivant traite l'intégration de cette chaîne avec le système sous test, puis la recette globale de l'outil par l'ingénieur BMS.
